\chapter*{Resumen}
\addcontentsline{toc}{chapter}{Resumen}
Las rachas secas (dry spells) condicionan la disponibilidad de agua de forma que la precipitación acumulada no captura. Este trabajo caracteriza la distribución de duraciones de dry spells en Chile continental con el producto observacional grillado CR2MET v2.5 y la simulación regional CNRM-ALADIN64 (dominio CHP12, \textasciitilde12 km), forzada por CNRM-ESM2-1, para 1980--2014 y para 2040--2074 bajo SSP5-8.5. Para comparar ambos productos de forma equitativa, se calibra el umbral de día húmedo de ALADIN (R₀ = 5,285 mm día$^{-1}$) de modo que su frecuencia media de días húmedos iguale la de CR2MET con 1 mm. ALADIN reproduce el gradiente norte--sur de persistencia, pero subestima la duración de las rachas, sobre todo en la cola: la duración media es 8,2 días menor y el percentil 99, 46,3 días menor que en CR2MET, con diferencias significativas en el 65 \% y el 43 \% de los píxeles, respectivamente. Hacia mediados de siglo, ALADIN proyecta menos rachas secas, pero más largas: el número de rachas de al menos 20 días iniciadas entre marzo y noviembre aumenta entre 12 \% y 115 \% según la región, con aumentos significativos en Coquimbo (RR = 1,12) y La Araucanía (RR = 1,72), y las distribuciones de Maipo, Maule y Biobío extienden su cola. La evaluación histórica es robusta al método de regrillado, y la dirección de la señal futura se mantiene con otro umbral, donde hay rachas suficientes para evaluarla, y con el escenario SSP3-7.0. El secado proyectado no es solo una reducción de la precipitación total, sino una reorganización temporal hacia períodos secos más persistentes; evaluar modelos para sequía requiere, por tanto, comparar la distribución completa de duraciones tras controlar la frecuencia de días húmedos.

\textbf{Palabras clave:} rachas secas, distribución de duraciones, modelo climático regional, CR2MET, ALADIN, Chile, cambio climático.

\chapter*{Abstract}
\addcontentsline{toc}{chapter}{Abstract}
\pendiente{traducir el resumen al inglés cuando esté cerrado.}

\chapter*{Agradecimientos}
\addcontentsline{toc}{chapter}{Agradecimientos}
\pendiente{redactar.}

\chapter*{Acrónimos y símbolos}
\addcontentsline{toc}{chapter}{Acrónimos y símbolos}
{\def\LTcaptype{none} % do not increment counter
\begin{longtable}[]{@{}
  >{\raggedright\arraybackslash}p{(\linewidth - 2\tabcolsep) * \real{0.5000}}
  >{\raggedright\arraybackslash}p{(\linewidth - 2\tabcolsep) * \real{0.5000}}@{}}
\toprule\noalign{}
\begin{minipage}[b]{\linewidth}\raggedright
Término
\end{minipage} & \begin{minipage}[b]{\linewidth}\raggedright
Significado
\end{minipage} \\
\midrule\noalign{}
\endhead
\bottomrule\noalign{}
\endlastfoot
ALADIN & Modelo climático regional CNRM-ALADIN64 \\
CDD & Consecutive Dry Days: racha seca más larga del año (índice ETCCDI) \\
CHP12 & Dominio de la simulación regional sobre Chile, \textasciitilde12 km \\
CMIP6 & Coupled Model Intercomparison Project, fase 6 \\
CR2MET & Producto grillado de precipitación del Centro de Ciencia del Clima y la Resiliencia (CR2) \\
ETCCDI & Expert Team on Climate Change Detection and Indices \\
GCM / RCM & Modelo climático global / regional \\
PDF & Función de densidad de probabilidad; aquí, densidad de frecuencia de duraciones \\
RR(D) & Risk ratio: N\_fut(t ≥ D) / N\_hist(t ≥ D) \\
SSP & Shared Socioeconomic Pathway (escenario de emisiones) \\
R, R₀ & Umbral de día húmedo; R₀ es el umbral calibrado de ALADIN (5,285 mm día$^{-1}$) \\
t, D & Duración de una racha; umbral de duración en los risk ratios \\
f(t) & Densidad de frecuencia: rachas por bin dividido por el ancho del bin \\
t99, Δ & Percentil 99 de las duraciones; diferencia ALADIN − CR2MET \\
\end{longtable}
}

